\chapter*{Abstract}
Large Language Models produce fluent but factually incorrect responses at rates reported up to 40\% in open-domain question answering, and higher in specialised domains. They also cannot improve on their own because they operate statelessly, without verified knowledge accumulation across interactions. This research proposes CAEM (Confidence-Aware Episodic Memory with Self-Improvement), a system that progressively reduces the specific subclass of confident \textit{factuality} hallucinations in knowledge-based question answering: confabulations~\cite{farquhar2024semantic}, ungrounded factual claims, and claims that contradict retrieved evidence. CAEM integrates three core mechanisms. First, a verified episodic memory populated through a four-outcome storage decision (\textsc{Store}, \textsc{Deferred}, \textsc{Abstain}, \textsc{Discard}) driven by a nine-signal post-generation verifier organised in three complementary families (internal calibration, sample-set agreement, external grounding), with an ensemble-max contradiction probability as a separately-consumed hard veto rather than a weighted composite input. Second, a confidence-aware three-tier router that combines memory-match similarity with stored verifier confidence into a single dispatch score, and an independent safety override that forces retrieval-augmented generation whenever pre-routing confidence falls below a fixed floor. Third, a constrained self-improvement loop that, at each cycle boundary, retroactively re-verifies every stored episode under the updated model, reconsiders deferred entries, and fine-tunes under an $L_2$ anchor toward the base-model parameters as an EWC approximation; a hard MMLU retention guard rolls back the weights (but not the memory) when general capability erodes beyond a fixed floor. The system commits to a pre-registered target of at least a 40\% relative reduction in the pooled confident-error rate against a zero-shot Flan-T5-Large baseline on a six-benchmark factual-QA panel, with MMLU as the out-of-distribution retention control. The study also argues, under a data-purity analysis, that even an imperfect verifier can drive monotonic improvement across cycles when the training-pool purity threshold is set above the storage threshold.


\vspace{0.5cm}

\noindent\textbf{Keywords:} Large Language Models, Hallucination Reduction, Episodic Memory, Confidence Estimation, Self-Improvement, Verification Systems, Continual Learning